\documentclass[11pt]{article}
% DRAFT seed (2026-10-13). Numbers come from numbers.tex (habitat/manuscript_numbers.py); never type a result here.
\usepackage[T1]{fontenc}
\usepackage[utf8]{inputenc}
\usepackage[margin=25mm]{geometry}
\usepackage{authblk}
\usepackage[numbers,sort&compress]{natbib}
\usepackage{graphicx}
\usepackage{booktabs}
\usepackage{xcolor}
\usepackage{url}
\usepackage[hidelinks]{hyperref}
\newcommand{\todo}[1]{\textcolor{red}{[TODO: #1]}}
% generated by habitat/manuscript_numbers.py from the result files; do not edit by hand
\newcommand{\LedgerDigest}{53413134410685ff}
\newcommand{\GitHead}{b00d857}
\newcommand{\NRecords}{174}
\newcommand{\NAssumed}{80}
\newcommand{\NSpecified}{51}
\newcommand{\NDerived}{42}
\newcommand{\FigTwoMaxErrIn}{0.092}
\newcommand{\FigTwoDims}{22}
\newcommand{\BaseCells}{255{,}123}
\newcommand{\BaseVolume}{63.1}
\newcommand{\BaseHeatErr}{+0.14}
\newcommand{\BaseCOTwoErr}{-0.57}
\newcommand{\BaseAgeErr}{-0.16}
\newcommand{\BaseReturnCOTwo}{88}
\newcommand{\GciHFine}{48.7}
\newcommand{\GciHMed}{62.8}
\newcommand{\GciHCoarse}{84.6}
\newcommand{\GciDomainMax}{2.7}
\newcommand{\GciLocalMax}{13}
\newcommand{\GciRoomT}{2.9}
\newcommand{\SstLo}{0.07}
\newcommand{\SstHi}{0.18}
\newcommand{\KeLo}{0.04}
\newcommand{\KeHi}{0.12}
\newcommand{\OffgasMixed}{0.8}
\newcommand{\OffgasCanopyNFT}{1.1}
\newcommand{\LoadMax}{450}
\newcommand{\LightMaxPPF}{145}
\newcommand{\LightPowerNeeded}{622}
\newcommand{\LightEff}{3}
\newcommand{\MecLo}{200}
\newcommand{\MecHi}{500}
\newcommand{\WallArea}{9.33}
\newcommand{\SthreeSixCrew}{2.1}
\newcommand{\SthreeTwoCrew}{6.3}
\newcommand{\OutdoorCOTwo}{426.1}
\newcommand{\SthreeVolume}{63.1}
\newcommand{\AisleMin}{0.89}
\newcommand{\WalkArea}{6.5}
\newcommand{\FloorArea}{25.2}
\newcommand{\ModelDominates}{\todo{pending: k-epsilon baseline}}
\newcommand{\ModelQuantities}{\todo{pending: k-epsilon baseline}}


\title{A sourced parametric model of an analog-habitat food production module: CAD, airflow CFD and a shared VR
review built from one ledger \todo{title to be agreed}}

\author[1]{Richard Barker\thanks{ORCID 0000-0001-5681-9857. \todo{corresponding-author email and any further affiliation to confirm}}}
\author[2]{John M. Peters\thanks{\todo{confirm full name as it should be printed (the ICES-2026-251 byline gives ``John M. Peters'') and ORCID}}}
\author[1]{D. Marshall Porterfield\thanks{\todo{confirm name form and ORCID}}}
\affil[1]{Agricultural \& Biological Engineering, Purdue University, West Lafayette, IN 47907, USA}
\affil[2]{Space \& Earth Analogs Research Chapter, Purdue University, West Lafayette, IN 47907, USA}
\date{Draft, ledger \LedgerDigest, code \GitHead}

\begin{document}
\maketitle

\begin{abstract}
Analog habitats are one place where food-production and biomanufacturing systems for crewed missions can be
operated together, yet the design data for such facilities are usually scattered across papers, drawings and renders. We
built a parametric model of the ASTRO-USA Food Production Module (FPM), a converted shipping container described
in \citet{peters2026astro}, in which one ledger of \NRecords{} values (\NSpecified{} quoted from sources,
\NDerived{} derived from them, \NAssumed{} declared assumptions) generates engineering CAD, airflow simulations and a shared VR scene. The
OpenFOAM simulations solve buoyant, time-averaged turbulent flow with heat, CO\textsubscript{2}, water vapour and age
of air, and every run must close its balances within 1\,\% before it is used. A ventilated-room benchmark (IEA
Annex 20) favoured the standard $k$--$\varepsilon$ model over SST $k$--$\omega$ (RMS velocity errors \KeLo--\KeHi{}
against \SstLo--\SstHi{} of the inlet velocity). With a placeholder HVAC unit, room-wide quantities are robust to mesh
and turbulence model (within \GciDomainMax{} and \ModelDomainMax\,\%), single-point values are not (up to
\GciLocalMax{} and \ModelLocalMax\,\%). Halving the airflow lowers the air-change efficiency from \BaseACE{} to
\SfiveaACE; CO\textsubscript{2} off-gas from placeholder bioreactor payloads is about \OffgasMixed~ppm against
\BaseReturnCOTwo~ppm from two crew; an airtight module with the HVAC off reaches the OSHA CO\textsubscript{2} limit
in \SthreeSixCrew~h with six crew; and the stated \LoadMax~W power budget, at the best LED efficacy achieved to date, cannot light the
hydroponic wall into the range of NASA's lettuce crop model. The model also yields a sensor position that stays representative across all
operating conditions run, and a list of the facts that would turn its assumptions into a design tool.
\end{abstract}

\section{Introduction}
Long-duration crewed missions are expected to rely in part on growing food \citep{wheeler2017agriculture}, and ground
analogs are where such systems can be operated with people in the loop before flight. Examples range from the EDEN ISS greenhouse operated in Antarctica
\citep{zabel2020edeniss} to crewed bioregenerative life-support experiments \citep{fu2016lunarpalace}; the history of crop research for space
agriculture is reviewed by \citet{wheeler2017agriculture}. In these facilities the air is shared:
crew, crops, lights and any biomanufacturing hardware all add heat, CO\textsubscript{2} and water vapour to one
ventilated volume, so where things are placed, how the air is moved and where it is measured are design decisions.

ASTRO-USA is an analog habitat being developed at Purdue University \citep{peters2026astro}. Its first physical module, the
Food Production Module, is a converted shipping container holding a vertical nutrient-film (NFT) hydroponic wall of
\NTowers{} towers with \NSitesPerTower{} sites each (\NSites{} plant sites), smaller hydroponic systems, a
bay for bioreactor payloads, a battery bank and a work space, with an HVAC unit in
an attic at one end. The paper gives a dimensioned floor plan and an electrical budget of \LoadMax~W maximum and
\LoadAvg~W average \citep{peters2026astro}; much else needed to predict the module's air (the HVAC unit, the
lighting, the envelope and the exact payloads) is not yet specified.

Computational fluid dynamics is an established tool for room air distribution \citep{nielsen2015fifty}, but its
predictions depend on the turbulence model \citep{chen1995kepsilon} and the mesh, so they need validation and an
uncertainty estimate \citep{celik2008gci}. They also depend on geometry and loads that, for a facility still being
designed, change often and are held in different places by different people. A model that is only as good as its
least-documented input is hard to trust or update.

Here we build the FPM from a single sourced ledger in which every value is either quoted from a source or declared as
an assumption, and generate from it engineering CAD, CFD cases and a VR scene for shared design review, so that a
corrected input propagates to every output. We validate the CFD set-up against a ventilated-room benchmark, quantify
its mesh and turbulence-model uncertainty, and use it to examine occupancy, airflow, biomanufacturing off-gas siting,
lighting power, HVAC failure and sensor placement. Because many inputs are still assumptions, the results are a
screening of the design space and a list of the facts the team can supply to sharpen it, not predictions of the
built module.

\section{Methods}
\subsection{Sources and the ledger}
Every dimension and load is a ledger record carrying its value, unit and either a source with a verbatim quote or a
declared assumption with its reason; the build refuses anything else. Of the \NRecords{} records, \NSpecified{} are
quoted from sources, \NDerived{} are derived from them (for example by vector extraction) and \NAssumed{} are
assumptions, most of which only the ASTRO-USA team can replace. The floor plan was extracted from the vector
geometry of Figure~2 of \citet{peters2026astro}: all \FigTwoDims{} printed dimensions agree with their own dimension
lines within \FigTwoMaxErrIn{}~in. Container dimensions come from a manufacturer's specification \citep{seacube40hc}
and crew metabolic loads from BVAD Table~3-26 \citep{bvad2022}.

\subsection{Geometry tiers}
A parts registry turns the ledger into engineering CAD (build123d; STEP, STL, GLB and dimensioned DXF drawings), into
CFD geometry with named patches and source zones, and into a procedural VR scene at three levels of detail with
per-node metadata, colliders and a walkable navigation mesh. The bioreactor payloads are generic placeholders, represented
by their outer envelopes, gas ports and loads only.

\subsection{Airflow CFD}
OpenFOAM \citep{openfoam2606} in a digest-pinned container. The living space (\BaseVolume{}~m\textsuperscript{3}
of air) is meshed with snappyHexMesh; the attic HVAC unit is represented by a supply and a return grille with equal
mass flow. CO\textsubscript{2}, water vapour and the age of air are transported as scalars relative to the supply
air. Steady RANS runs did not settle (a cold supply jet in a heated room), so each case is solved as a transient
(buoyantPimpleFoam) over three room flush times ($V/Q$) and averaged over the last two. Two cases needed longer:
the payload-siting case was continued for three more flush times with a fresh averaging window, and the half-airflow
case, which was still settling, was run uninterrupted for six flush times and averaged over the last three (a
restarted extension of it was discarded because the restart itself disturbed the solution; see Discussion). Every run must close its mass, heat,
CO\textsubscript{2}, H\textsubscript{2}O and age-of-air balances within 1\,\% and be statistically stationary before
any result is used. Discretisation uncertainty follows \citet{celik2008gci} on three meshes (effective cell sizes
\GciHFine, \GciHMed{} and \GciHCoarse~mm); a GCI is quoted only where the mesh-to-mesh differences exceed the
time-averaging noise and the observed order is plausible.

\subsection{Validation and turbulence model}
The IEA Annex 20 two-dimensional isothermal test case \citep{nielsen1990annex20} was solved with the habitat cases'
wall functions and schemes on three meshes and compared with laser-Doppler measurements \citep{aalborg2006lda}.
The SST $k$--$\omega$ model \citep{menter1994} had RMS errors of \SstLo--\SstHi{} of the inlet velocity and the
standard $k$--$\varepsilon$ model \citep{launder1974} \KeLo--\KeHi; the scenario runs therefore use
$k$--$\varepsilon$. The buoyant variant of the benchmark has no measurements.

\section{Results}
\subsection{Baseline}
The time-averaged baseline (\BaseCells{} cells) closes its balances with errors of \BaseHeatErr\,\% (heat),
\BaseCOTwoErr\,\% (CO\textsubscript{2}) and \BaseAgeErr\,\% (age of air). Across the three meshes, room-wide
means change by at most \GciDomainMax\,\% and local values by up to \GciLocalMax\,\%; the room-mean temperature rise
has a GCI of \GciRoomT\,\% (observed order \GciRoomP, at the lower edge of the plausible range, so a
conservative figure). Repeating the baseline with $k$--$\varepsilon$ changes room-wide means by at most
\ModelDomainMax\,\%, local values by up to \ModelLocalMax\,\% and air speeds by up to \ModelSpeedMax\,\%: the
room-scale picture holds under both mesh and turbulence model, while single-point values remain uncertain under both.

\subsection{Occupancy and airflow}
With six crew at once the room-mean CO\textsubscript{2} rise grows $\times$\StwoCOTwoRatio{} (to \StwoCOTwo~ppm above the supply air) and the temperature rise $\times$\StwoTRatio. Halving the airflow raises the temperature rise $\times$\SfiveaTRatio, CO\textsubscript{2} $\times$\SfiveaCOTwoRatio{} and the age of air $\times$\SfiveaAgeRatio{}, more than the $\times$2 of a perfectly mixed room: mixing deteriorates at low flow. Expressed as the air-change
efficiency (the nominal time constant divided by twice the room-mean age of air, 0.5 for complete mixing;
\citealp{sandberg1983moments}), the baseline is at \BaseACE{} and half airflow at \SfiveaACE. The double-airflow case did not reach a statistically stationary state within its run and is not reported.

\subsection{Bioreactor off-gas siting}
The placeholder bioreactor rack's off-gas adds about \OffgasMixed~ppm of CO\textsubscript{2} when mixed through the room, against
about \BaseReturnCOTwo~ppm from two working crew; piped to the hydroponic wall it raises the canopy zone by
\OffgasCanopyNFT~ppm ($k$--$\varepsilon$ run).

\subsection{Lighting power}
If the stated \LoadMax~W maximum draw includes the grow lights, the \WallArea~m\textsuperscript{2} NFT wall cannot exceed
\LightMaxPPF~$\mu$mol\,m$^{-2}$\,s$^{-1}$ even with every watt going to light at the best achieved fixture efficacy
(\LightEff~$\mu$mol\,J$^{-1}$) \citep{kusuma2020led}, below the \MecLo--\MecHi~$\mu$mol\,m$^{-2}$\,s$^{-1}$
range of BVAD's lettuce crop model \citep{bvad2022}; \LightPowerNeeded~W would be needed for lights alone.

\subsection{HVAC failure}
With the HVAC off and the door shut, CO\textsubscript{2} in the \SthreeVolume~m\textsuperscript{3} of air rises
from the outdoor level (\OutdoorCOTwo~ppm) \citep{noaa2026co2} to the OSHA permissible exposure limit
\citep{osha1910} in \SthreeSixCrew~h with six crew and \SthreeTwoCrew~h with two if the container is airtight.

\subsection{Layout and sensors}
The narrowest clear aisle is \AisleMin~m; keeping clear of equipment leaves \WalkArea{} of the
\FloorArea~m\textsuperscript{2} floor walkable. Scoring every wall position against the occupied-zone mean over \SensorRuns{} passing $k$--$\varepsilon$ runs, the most representative single sensor sits on \SensorWhere{} at \SensorZ~m, within \SensorT~K and \SensorCOTwo\,\% (CO\textsubscript{2} excess) in every run and \SensorClear~m from the nearest growing system.

\section{Discussion}
\paragraph{What the model can say now.} With the HVAC unit, lighting and envelope still unknown, the model is a
screening tool. Its most robust outputs are room-scale: room-mean temperature rise, CO\textsubscript{2} and age of
air changed by a few per cent between meshes and between turbulence models, whereas values at single points (a
probe beside a crewmember, a position along the aisle) changed by up to \GciLocalMax--\ModelLocalMax\,\%. Decisions
that depend on room-scale balances (how much CO\textsubscript{2} a source adds, how ventilation rate trades against
temperature) can therefore be explored now; decisions that hinge on the air at a specific point need finer meshes
and, above all, measurements.

\paragraph{Ventilation.} Halving the airflow more than doubled the room's CO\textsubscript{2} and temperature rise
and reduced the air-change efficiency from \BaseACE{} to \SfiveaACE, so the room mixes less well at low flow, not
only less often. The airflow of the attic unit, and whether it brings in outdoor air at all, therefore matter more
than any placement question studied here. If the unit only recirculates room air, CO\textsubscript{2} accumulates as
in the HVAC-off case: \SthreeSixCrew~h to the OSHA limit with six crew in an airtight module. The limit is an
eight-hour time-weighted average, and real leakage would slow the rise, so this is a planning bound rather than a
hazard prediction; a door-shut CO\textsubscript{2} decay test would measure the leakage that decides it.

\paragraph{Bioreactor payloads.} At the placeholder off-gas rate, routing the bioreactors' CO\textsubscript{2}
to the hydroponic wall raises the canopy by about \OffgasCanopyNFT~ppm, which is small beside the crew's contribution.
At these rates the payloads are neither a CO\textsubscript{2} hazard nor a useful CO\textsubscript{2} supply for the
plants; that conclusion would change only if their real fluxes are much larger than the planning values used here.

\paragraph{Lighting and plants.} If the \LoadMax~W budget includes the grow lights, the hydroponic wall cannot reach
the light levels that NASA's lettuce crop model assumes with fixtures at the best efficacy achieved to date. Either the lighting sits outside
the stated budget or the crop will be strongly light-limited; which of the two holds changes the expected yield and is
the main reason plant gas exchange is not yet coupled into the CFD. Until it is, the plants appear only as LED heat,
which is why the sensor ranking excludes positions next to them.

\paragraph{Sensor placement.} A single wall sensor on \SensorWhere{} at \SensorZ~m stayed within \SensorT~K and
\SensorCOTwo\,\% of the occupied-zone mean across \SensorRuns{} operating conditions. That makes it a reasonable first
position for the module's monitoring and, once logging, a source of the measurements the model needs for validation.

\paragraph{Lessons for modelling small ventilated habitats.} Three practical points emerged. Steady RANS did not
settle for a cold supply jet in a heated room, so time-averaged transients with explicit balance and stationarity
gates were needed. With fixed mass flows at both grilles the room behaves as a sealed volume whose pressure rises as
it warms; restarting such a run after a large pressure drift injected energy into the solution (the restarted
half-airflow case showed a step in room-mean temperature at the restart), so long cases are better run uninterrupted. And mesh-uncertainty estimates
from time-averaged fields are fragile when mesh-to-mesh differences approach the averaging noise; reporting the
spread alongside the GCI, and quoting a GCI only where both tests pass, avoids over-reading them.

\paragraph{Limitations.} The HVAC unit, its supply temperature and outdoor-air fraction are placeholders; the walls
are adiabatic; meshes are coarse; the validation is isothermal only; plants and fungi exchange no gas; payloads are
modelled as envelopes with planning-value fluxes; and the double-airflow case did not reach a stationary state within
its run and is not reported. Each of these is a ledger entry or a TEAM\_INPUTS item, so replacing it with the team's
data re-runs the same pipeline.

\section*{Data and code availability}
Code, ledger and results: \todo{repository URL once public, and Zenodo DOI}. The ICES-2026-251 paper and slides
are not redistributed.

\section*{Author contributions}
\todo{draft for the authors to confirm or replace (CRediT):} R.B.: conceptualisation, methodology, software,
validation, formal analysis, writing (original draft). J.M.P.: \todo{for example ASTRO-USA FPM design and source
data; writing (review and editing)}. D.M.P.: \todo{for example supervision; writing (review and editing)}.

\section*{Use of AI tools}
\todo{draft for the authors to confirm against the target journal's policy:} The model code, simulations, analyses
and a first draft of this text were produced with an AI coding assistant (Claude, Anthropic) under the direction of
R.B. Every source value, citation and reported number was checked against its primary source or the generated result
files, as documented in the repository; the authors take responsibility for the content.

\section*{Competing interests}
\todo{to be declared by the authors.}

\bibliographystyle{plainnat}
\bibliography{refs}
\end{document}
